%%
%% Elsevier CAS double-column draft organized for the algorithm ablation study.
%% Replace the author and affiliation metadata below with the final submission info.
%%
\documentclass[a4paper,fleqn]{cas-dc}

\usepackage[numbers]{natbib}
\usepackage{amsmath}
\usepackage{capt-of}
\usepackage{needspace}
\flushbottom

\begin{document}
\let\WriteBookmarks\relax
\def\floatpagepagefraction{1}
\def\textpagefraction{.001}

\shorttitle{Algorithm Ablation of Separated-Skeleton MATD3}
\shortauthors{Author et al.}

\title[mode=title]{Algorithm Ablation of a Separated-Skeleton MATD3 for Terrain-Aware Multi-UAV Coordination}

\author[1]{Author Name}
\cormark[1]
\ead{author@example.com}

\affiliation[1]{organization={Your Institution},
                city={Your City},
                country={Your Country}}

\cortext[cor1]{Corresponding author.}

\newcommand{\ablationplotdir}{../matd3/ablation_experiments/multi_seed_groupB_20260331_220752/plots}

\begin{abstract}
This paper organizes the current codebase into a focused algorithm-ablation study for terrain-aware multi-UAV coordination. The mainline method is a MATD3 variant with a separated update skeleton over raw and corrected action semantics, together with a separated actor objective. We compare it against three MATD3 counterparts: the same skeleton with a unified actor objective, a hybrid actor objective, and a standard single-total-$Q$ MATD3 baseline. Three MADDPG variants are included as cross-family references. All methods are evaluated under a strict protocol with fixed training positions, a shared terrain seed, disabled curriculum progression, and three random seeds. On the Group-B random-obstacle benchmark, the mainline reaches a tail-100 team success rate of 58.33\%, compared with 43.33\% for the unified-objective ablation and 27.00\% for the MATD3 baseline. It also achieves the highest tail-100 reward (243720.59) and the lowest collision count among the MATD3 variants (22.20). In the shared post-training evaluation, the mainline is the only MATD3 variant that retains non-zero mean team success. These results support the conclusion that both the separated skeleton and the separated actor objective contribute materially to coordinated learning performance.
\end{abstract}

\begin{highlights}
\item A strict multi-seed ablation protocol isolates the critic skeleton from the actor-objective design.
\item The separated actor objective improves tail-100 team success from 43.33\% to 58.33\% under the same MATD3 skeleton.
\item The separated-skeleton MATD3 mainline outperforms the MATD3 baseline and all MADDPG references on the studied benchmark.
\end{highlights}

\begin{keywords}
multi-agent reinforcement learning \sep MATD3 \sep ablation study \sep terrain-aware coordination \sep actor objective \sep dual-semantic critic
\end{keywords}

\maketitle

\section{Introduction}

Terrain-aware multi-UAV coordination is difficult because the learning signal must jointly capture collision avoidance, target reaching, and coordinated team completion under complex 3D geometry. In practice, strong performance often comes from several interacting implementation choices at once, which makes it hard to determine which module is actually responsible for the gain. For this reason, the current codebase benefits from being reorganized as a paper-style ablation study instead of only a best-method report.

The present implementation already defines a clear mainline method: a MATD3 framework with separated raw/corrected action semantics, a dual-component value decomposition, and a separated actor objective. Around this mainline, the ablation script constructs matched variants that keep the outer training loop, correction interface, and environment setup fixed while changing only one algorithmic factor at a time. This makes it possible to answer two central questions. First, does the separated actor objective help when the separated update skeleton is kept unchanged? Second, does the separated skeleton itself still help when compared with a standard MATD3 baseline?

This paper is therefore organized around the ablation experiment implemented in \texttt{ablation\_dual\_q\_separated\_gradient.py}. Under a strict protocol with fixed positions, the same terrain seed, disabled curriculum, and three random seeds, the mainline method achieves the best training reward, the highest team success rate, and the lowest collision count among the MATD3 variants. The results show that the separated actor objective yields a clear gain over the unified alternative, while the dual-semantic skeleton also improves coordination relative to the baseline single-total-$Q$ design.

\section{Algorithm Variants Under Study}

\subsection{Task-facing policy structure}

The ablation is conducted on the terrain-aware multi-UAV coordination task used by the current training code. Each policy outputs the same 7-D action representation used throughout the implementation, and all compared methods keep the same correction interface enabled through \texttt{USE\_TF\_POTENTIAL\_FIELD=1}. This design choice is important: it prevents the ablation from being confounded by switching the action-correction pipeline on and off. The comparison is therefore centered on the learning architecture rather than on the presence or absence of the correction module.

\subsection{Separated skeleton and actor objectives}

The mainline method decomposes the critic signal into two semantic parts:
\begin{equation}
Q_{\text{tot}}(s, a^{\text{raw}}, a^{\text{corr}})
=
Q_{\text{head}}(s, a^{\text{raw}})
+
Q_{\text{tail}}(s, a^{\text{corr}}),
\label{eq:dual-q}
\end{equation}
where $a^{\text{raw}}$ denotes the raw action semantics and $a^{\text{corr}}$ denotes the corrected or executed action semantics. The purpose of this design is to let the learning signal distinguish the policy intent from the action after correction.

The core ablation then studies three actor objectives on top of this skeleton:
\begin{align}
J_{\text{sep}}   &= \mathbb{E}\!\left[Q_{\text{head}} + Q_{\text{tail}}\right], \\
J_{\text{uni}}   &= \mathbb{E}\!\left[Q_{\text{tot}}\right], \\
J_{\text{hyb}}   &= \alpha J_{\text{sep}} + (1-\alpha)J_{\text{uni}},
\quad \alpha = 0.8.
\end{align}
Here, $J_{\text{sep}}$ is the default mainline objective, $J_{\text{uni}}$ is the matched ablation that collapses the actor update into a unified total-$Q$ target, and $J_{\text{hyb}}$ is a supplemental variant that mixes the two. In addition, a standard MATD3 baseline removes the separated decomposition and uses a conventional single-total-$Q$ critic per twin critic branch.

\subsection{Research questions}

The paper is organized around two direct ablation questions.

\begin{enumerate}
\item Does the separated actor objective improve performance when the separated MATD3 skeleton is fixed? This is answered by comparing M3-Sep with M3-Uni.
\item Does the separated skeleton improve performance over a standard MATD3 baseline? This is answered by comparing M3-Uni with M3-Base and by comparing M3-Sep with M3-Base as an end-to-end reference.
\end{enumerate}

\section{Experimental Protocol}

\subsection{Strict controlled setting}

All reported results follow the \texttt{strict\_ablation} configuration of the codebase. The training schedule uses 400 episodes, batch size 1024, weighted reward enabled, scenario seed 88, and the fixed position file \texttt{strict\_ablation\_seed88\_groupB.json}. Group B corresponds to the random-obstacle setting. Environment isolation is strict: curriculum unlocking is disabled, terrain progression is disabled, and all compared methods are trained on the same fixed training layout under the same random seeds $\{936487, 101, 202\}$.

The shared post-training evaluation uses the best checkpoint selected by team success rate, a shared train-match test mode (the current implementation canonicalizes the legacy \texttt{heldout\_shared} alias into this shared test setting), 20 evaluation episodes, and post-evaluation seed 10088. Some variants do not appear in the post-training table because they never produced a qualified checkpoint with team success during training; in those cases the script explicitly skips post-evaluation.

The action-correction force ratio (FR) is treated as part of the checkpoint semantics rather than as a single fixed evaluation constant. During training, all methods share the same FR schedule, initialized at 0.50 and gradually annealed to 0.32. This decreasing schedule was selected from preliminary tuning as the most stable setting used in the current experiments. The larger early FR provides stronger APF-mediated guidance when the policies are still poorly shaped, reducing unsafe exploration and helping the agents discover feasible motion. As training progresses, the FR is reduced so that the learned policy carries more of the control burden and does not overfit to a strongly corrected execution interface. During post-training evaluation, the FR schedule is disabled because no further training progression occurs, but the scalar FR is restored from the selected checkpoint record: \texttt{best\_by\_team\_sr} uses \texttt{best\_team\_sr\_force\_ratio}, \texttt{best} uses \texttt{best\_episode\_force\_ratio}, and \texttt{final}/\texttt{checkpoint} use the final recorded FR. Therefore, the reported post-training results use checkpoint-consistent deployment FR values, not a globally fixed FR. Any separate deployment-control sweep that intentionally fixes FR, for example at 0.50, should be reported as a distinct controlled protocol rather than merged with the checkpoint-consistent evaluation.

\subsection{Metrics}

Two groups of metrics are reported.

\begin{enumerate}
\item Training metrics: tail-100 average reward, tail-100 team success rate, tail-100 collision count, and tail-100 average minimum clearance.
\item Post-training metrics: team success rate, average reward, average collision count, minimum clearance, team total path length, and success-arrival step when at least one successful episode exists.
\end{enumerate}

The training metrics are aggregated over the last 100 episodes of each seed, and the paper reports mean and standard deviation across the three seeds.

\section{Results}

\subsection{Training-phase ablation results}

Table~\ref{tab:training-results} summarizes the multi-seed training statistics. The mainline M3-Sep achieves the highest tail-100 reward and the highest team success rate among all variants. Under the same separated update skeleton, replacing the separated actor objective with the unified objective reduces success from 58.33\% to 43.33\%, lowers reward by 83096.66, and increases collisions from 22.20 to 28.95. This directly supports the claim that the actor-objective design matters even when the skeleton is kept fixed.

The MATD3 baseline M3-Base further clarifies the role of the separated skeleton. Compared with M3-Base, the mainline improves tail-100 team success by 31.33 percentage points and reduces collisions by 31.57. Even the matched unified-objective variant M3-Uni is clearly stronger than the standard baseline in team success and collision control, which indicates that the structural separation itself contributes useful learning signal.

The supplemental hybrid variant M3-H80 sits between M3-Sep and M3-Uni in reward and success, but it is notably worse in collision count than both. This suggests that reintroducing a unified total-$Q$ term may partly recover reward optimization while weakening the coordination benefits of the pure separated objective.

\begin{center}
\begin{minipage}{\columnwidth}
\captionof{table}{Compact training summary on Group-B random obstacles. Values are mean $\pm$ standard deviation over three seeds, computed from the last 100 training episodes.}
\label{tab:training-results}
\centering
\scriptsize
\setlength{\tabcolsep}{2pt}
\resizebox{\columnwidth}{!}{%
\begin{tabular}{lcccc}
\toprule
Method & Reward & Succ. (\%) & Coll. & Clr. (m) \\
\midrule
M3-Sep  & \textbf{243720.59 $\pm$ 9842.66} & \textbf{58.33 $\pm$ 2.49} & \textbf{22.20 $\pm$ 3.31} & 15.56 $\pm$ 0.24 \\
M3-Uni  & 160623.93 $\pm$ 64128.87 & 43.33 $\pm$ 17.21 & 28.95 $\pm$ 17.27 & 16.13 $\pm$ 0.45 \\
M3-H80  & 187887.84 $\pm$ 35144.30 & 44.67 $\pm$ 6.02 & 41.83 $\pm$ 12.04 & 16.13 $\pm$ 0.21 \\
M3-Base & 160878.30 $\pm$ 52710.19 & 27.00 $\pm$ 8.60 & 53.77 $\pm$ 21.10 & 17.49 $\pm$ 0.52 \\
DG-Sep  & 3961.03 $\pm$ 14098.99 & 0.00 $\pm$ 0.00 & 59.57 $\pm$ 24.89 & 25.75 $\pm$ 3.08 \\
DG-Uni  & 15657.68 $\pm$ 42580.61 & 3.00 $\pm$ 4.24 & 59.05 $\pm$ 36.84 & 24.07 $\pm$ 3.11 \\
DG-Base & -165960.97 $\pm$ 106336.45 & 0.00 $\pm$ 0.00 & 301.22 $\pm$ 202.07 & 41.51 $\pm$ 7.59 \\
\bottomrule
\end{tabular}%
}
\end{minipage}
\end{center}

\subsection{Cross-family reference comparison}

The MADDPG reference family does not approach the MATD3 performance level in this benchmark. DG-Sep and DG-Base never achieve non-zero mean team success, while DG-Uni reaches only 3.00\% tail-100 success. Their reward is also far below the MATD3 family. These rows are therefore best interpreted as boundary references that confirm the difficulty of the benchmark and the importance of the MATD3-style stabilizing machinery, rather than as the primary basis for causal attribution inside the mainline method.

\subsection{Post-training shared evaluation}

The shared post-training evaluation is harder than the training metric itself, and the absolute team-success numbers remain low for all methods. Even so, the relative ranking is informative. As shown in Table~\ref{tab:post-eval-results}, M3-Sep is the only MATD3 variant with non-zero mean team success (1.67\%). Both M3-Uni and M3-Base drop to zero. The baseline also exhibits much longer team path length (1235.75 m), indicating weaker coordinated completion behavior even when a checkpoint is available for evaluation.

At the same time, the post-training metrics reveal that generalization is still a major challenge. M3-Uni obtains fewer collisions and a less negative average reward than M3-Sep in the shared evaluation, yet it still fails to complete the team task successfully. This suggests that conservative motion alone is not sufficient; what matters is whether the policy can preserve coordinated progress under distribution shift. In that sense, the mainline method remains the only variant that transfers any successful team completion to the shared post-training test.

\begin{center}
\begin{minipage}{\columnwidth}
\captionof{table}{Compact post-training shared evaluation. Only variants that produced an eligible checkpoint through training success are reported; evaluation restores the FR associated with the selected checkpoint rather than using one global fixed FR.}
\label{tab:post-eval-results}
\centering
\scriptsize
\setlength{\tabcolsep}{2pt}
\resizebox{\columnwidth}{!}{%
\begin{tabular}{lccccc}
\toprule
Method & Succ. (\%) & Reward & Coll. & Clr. (m) & Path (m) \\
\midrule
M3-Sep  & \textbf{1.67 $\pm$ 2.36} & -167312.61 $\pm$ 50566.99 & 319.40 $\pm$ 204.39 & 8.18 $\pm$ 1.21 & 890.12 $\pm$ 2.96 \\
M3-Uni  & 0.00 $\pm$ 0.00 & \textbf{-88225.12 $\pm$ 54895.67} & \textbf{53.20 $\pm$ 40.25} & 7.58 $\pm$ 1.17 & \textbf{888.25 $\pm$ 92.63} \\
M3-Base & 0.00 $\pm$ 0.00 & -155534.87 $\pm$ 42337.02 & 187.15 $\pm$ 73.07 & -7.37 $\pm$ 26.34 & 1235.75 $\pm$ 38.27 \\
DG-Uni  & 0.00 $\pm$ 0.00 & -79902.80 $\pm$ 0.00 & 37.25 $\pm$ 0.00 & 88.39 $\pm$ 0.00 & 2304.65 $\pm$ 0.00 \\
\bottomrule
\end{tabular}%
}
\end{minipage}
\end{center}

The supplementary training and post-evaluation plots are placed after Appendix~A.

\section{Discussion}

Three conclusions emerge from the ablation. First, the separated actor objective is not a cosmetic change. Under a fixed separated MATD3 skeleton, it yields the strongest reward, the highest success rate, and the fewest collisions. Second, the structural decomposition itself matters: the standard MATD3 baseline is substantially weaker in success rate and collision control. Third, the hybrid actor objective does not fully recover the behavior of the pure separated objective, which suggests that mixing unified and separated gradients may dilute the intended semantic split.

The post-training evaluation adds an important qualification. Although the mainline is the strongest overall variant, the shared test setting remains difficult and absolute success is still low. The paper should therefore avoid claiming that the problem is solved. A more appropriate interpretation is that the separated design improves coordination robustness, but further work is needed on generalization, checkpoint selection, and transfer calibration.

\section{Conclusion}

This paper reorganizes the existing implementation into a clean algorithm-ablation narrative for a terrain-aware multi-UAV MATD3 system. The evidence consistently favors the separated-skeleton MATD3 mainline, and the matched ablations show that both the actor-objective design and the dual-semantic critic structure matter. The resulting manuscript structure is suitable for extending into a full paper with additional related-work citations, broader benchmark coverage, and richer generalization analysis.

\appendix
\section{Exact Controlled Settings}

For reproducibility, the core controlled settings used in this manuscript are listed below.

\begin{enumerate}
\item Training mode: \texttt{strict\_ablation}.
\item Episodes: 400. Batch size: 1024. Weighted reward: enabled.
\item Scenario seed: 88. Position file: \texttt{strict\_ablation\_seed88\_groupB.json}.
\item Experiment group: Group B (random obstacles).
\item Training seeds: 936487, 101, and 202.
\item Post-training evaluation: shared train-match mode, 20 episodes, seed 10088, checkpoint selected by best team success rate.
\item Post-training FR protocol: training uses a shared scheduled FR, while evaluation disables the schedule and restores the scalar FR associated with the selected checkpoint (\texttt{best\_team\_sr\_force\_ratio}, \texttt{best\_episode\_force\_ratio}, or final recorded FR, depending on the checkpoint variant).
\item Core script: \texttt{ablation\_dual\_q\_separated\_gradient.py}.
\end{enumerate}

\begin{figure*}[t]
    \centering
    \includegraphics[width=.92\textwidth]{\ablationplotdir/matd3_trio_train_reward_success_20260407_073724.png}
    \caption{Training comparison among the three MATD3 variants most relevant to the main ablation claim. The mainline M3-Sep converges to higher reward and higher team success than both M3-Uni and M3-Base.}
    \label{fig:matd3-train}
\end{figure*}

\begin{figure*}[t]
    \centering
    \includegraphics[width=.95\textwidth]{\ablationplotdir/multi_seed_mean_ablation_comparison_20260407_073724.png}
    \caption{Multi-seed mean training curves for the full ablation set. The MATD3 family consistently dominates the MADDPG references on reward, success rate, and collision control in the studied setting.}
    \label{fig:full-train}
\end{figure*}

\begin{figure*}[t]
    \centering
    \includegraphics[width=.92\textwidth]{\ablationplotdir/matd3_trio_post_eval_reward_success_20260407_073724.png}
    \caption{Shared post-training evaluation for the three MATD3 variants. The mainline is the only variant that retains non-zero team success, although the entire setting remains challenging.}
    \label{fig:matd3-posteval}
\end{figure*}

\begin{figure*}[t]
    \centering
    \includegraphics[width=.95\textwidth]{\ablationplotdir/multi_seed_post_eval_summary_20260407_073724.png}
    \caption{Full post-training dashboard across seeds. The figure highlights a persistent gap between training performance and shared post-training performance, which leaves ample room for future generalization improvements.}
    \label{fig:full-posteval}
\end{figure*}

\clearpage
\bibliographystyle{cas-model2-names}
\bibliography{cas-refs}

\end{document}
